% Stage 02, WP S02-M2: strong approximation, the Brandt graph, the quotient graph.
% Labels are fixed by plan/staged/stage-02-tree-quotients.tex. Labels def:s02-tree, lem:s02-tree, def:s02-GU,
% def:s02-Xone, lem:s02-normone live in chapters/s02/1-tree.tex (WP S02-M1).
\subsection{Strong approximation and the Brandt graph}\label{ss:s02-quotient}

\paragraph{Conventions for this subsection.} $U$ is a $q$-admissible level (Def.~\ref{def:s02-GU}), so
$U=U_q\times U^{(q)}\subseteq\hat\cO^\times$ with $U_q=\iota_q^{-1}\GL_2(\Zq)$, and we put
\[U':=\Uzero{q}U=U\cap\Uzero{q},\qquad U'_q=\iota_q^{-1}\Iw_q,\qquad U'^{(q)}=U^{(q)}.\]
Matrices at $q$ are read through $\iota_q$; an element of $B_q^\times$ is regarded as the idele that is $1$ away from $q$.
For $\gamma\in B^\times$ (embedded diagonally in $\Bhat^\times$) we write $\gamma=\gamma_q\gamma^{(q)}$, where $\gamma_q\in B_q^\times$
and $\gamma^{(q)}$ is the idele with component $1$ at $q$ and $\gamma_\ell$ at $\ell\ne q$. Ideles supported at disjoint sets of
primes commute. Since $\Nrd$ is computed componentwise, $U^1=U^1_q\times U^{1,(q)}$ with $U^1_q=\iota_q^{-1}\SL_2(\Zq)$, and
likewise for $U'$. $\GU$ acts on $\cT_q$ by $\gamma\cdot P:=\iota_q(\gamma_q)P$.

\emph{Graphs.} All graphs here are \emph{bipartite multigraphs} $G=(V_0,V_1,E,o,t)$: sets $V_0,V_1,E$ and maps
$o:E\to V_0$, $t:E\to V_1$. A morphism $\phi:G\to H$ is a triple of maps $V_0\to V_0'$, $V_1\to V_1'$, $E\to E'$ with
$o\circ\phi=\phi\circ o$ and $t\circ\phi=\phi\circ t$. $G$ is \emph{connected} if $V_0\sqcup V_1\ne\emptyset$ and the
equivalence relation on $V_0\sqcup V_1$ generated by $o(e)\sim t(e)$ ($e\in E$) has one class; its classes, together with the
edges at their vertices, are the \emph{connected components}. The \emph{edges at} a vertex $w$ are the $e$ with $o(e)=w$
(if $w\in V_0$), resp.\ $t(e)=w$ (if $w\in V_1$).

\emph{Edge-set convention for $\cT_q$.} $V_j(\cT_q)$ is the set of type-$j$ vertices, $E(\cT_q)$ the set of (geometric,
unordered) edges $\{P,P'\}$; each has one end of each type (Def.~\ref{def:s02-tree}), and $o(e)$, $t(e)$ are its type-$0$
and type-$1$ ends. So $\cT_q$ is a bipartite multigraph in the above sense, and each edge carries its canonical orientation
type~$0\to$~type~$1$. (Serre's oriented edge set is $E\sqcup\bar E$; it is used for cochains in \S\ref{ss:s02-coh}, not here.)
$\GU$ acts by type-preserving automorphisms (Lemma~\ref{lem:s02-tree}(b)), hence commutes with $o$ and $t$. The
\emph{quotient graph} is
\[\GU\backslash\cT_q:=\big(\GU\backslash V_0(\cT_q),\ \GU\backslash V_1(\cT_q),\ \GU\backslash E(\cT_q),\ \bar o,\ \bar t\big),
\qquad \bar o(\GU e)=\GU o(e),\ \bar t(\GU e)=\GU t(e).\]

\emph{What is used about the tree.} Besides Def.~\ref{def:s02-tree}, the proofs below use only Lemma~\ref{lem:s02-tree}(b)--(d):
$\SL_2(\Qq)$ acts type-preservingly and transitively on $V_0$, on $V_1$ and on $E$, with stabilisers
\[S_{v_0}=\SL_2(\Zq),\qquad S_{v_1}=\SL_2(\Qq)\cap\alpha\GL_2(\Zq)\alpha^{-1},\qquad S_{e_0}=\SL_2(\Qq)\cap\Iw_q,\]
and $\alpha^{-1}\Iw_q\alpha\subseteq\GL_2(\Zq)$. This is exactly the coset model of the tree (Lean milestone L02.1). The fact that
$\cT_q$ is \emph{connected} (the tree property in Lemma~\ref{lem:s02-tree}(a), cited) is used at \emph{one} place only, marked
\textbf{[tree]} in the proof of Prop.~\ref{prop:s02-BG}.

\begin{theorem}[strong approximation for $B^1$ at $S=\{\infty,q\}$]\label{thm:s02-SA}
Since $q\nmid D$, $B_q\cong M_2(\Qq)$ is split and $B^1_q\cong\SL_2(\Qq)$ is non-compact. Hence $B^1$ (embedded diagonally) is
dense in the restricted product $\Bhat^{1,(q)}:=\prod'_{\ell\ne q}B^1_\ell$ (with respect to the $\cO^1_\ell$).
\smallskip

\noindent\emph{Consequence (the shape used; \texttt{if:s02-SA}).}\label{if:s02-SA} For every $q$-admissible $U$,
\[\Bhat^1=B^1\big(B^1_q\times U^{1,(q)}\big).\tag{SA$_U$}\]
\end{theorem}
\noindent\emph{Status.} The theorem is \stcited{}:
\lit{voight\_2021\_quaternion-algebras}{Main Thm.~28.5.3, p.~485}, applied to the global ring $R=\Z[1/q]$ with eligible set
$S=\{\infty,q\}$. In Voight's notation $\Bhat=\Bhat^S=\prod'_{v\notin S}B_v$ is the $S$-finite adele ring
(\lit{voight\_2021\_quaternion-algebras}{27.6.4, p.~469}; there, when $S$ is clear from the context, the superscript $S$ is dropped and replaced by hats), so his
$\Bhat^1$ is our $\Bhat^{1,(q)}$. The hypothesis ``$B$ is $S$-indefinite'' (\lit{voight\_2021\_quaternion-algebras}{Def.~28.5.1,
p.~485}: $S$ contains a place unramified in $B$) holds because $q\in S$ and $q\nmid D$.
Equivalently \lit{vigneras\_1980\_arithmetique-des-algebres-de-quaternions}{Ch.~III, Thm.~4.3, p.~81} ($H^1_KH^1_S$ dense in
$H^1_{\A}$ if $H^1_S$ is not compact; page-checked) and
\lit{platonov-rapinchuk\_1994\_algebraic-groups-and-number-theory\_2up}{\S7.4, Thm.~7.12, p.~427} (page-checked in the M8 litnote).
Kneser \lit{kneser\_1969\_lectures-on-galois-cohomology-of-classical-groups}{\S5.5, Sublemma~b, p.~94 (Tata pagination; p.~72 of the scan)} only states the result;
it is listed for orientation and not relied upon. The consequence (SA$_U$) is \stproved{} below from the theorem. It is
\lit{voight\_2021\_quaternion-algebras}{28.4.8, p.~484} (density implies $\Bhat^1\subseteq B^1\hat\cO^1$) and
\lit{voight\_2021\_quaternion-algebras}{Cor.~28.5.12, p.~486}; the shape (SA$_U$), with the factor $B^1_q$, is the projection away from $\infty$ of the \emph{second}
equality of (28.5.13) there (which multiplies in $B^1_S=B^1_\infty\times B^1_q$). Both are stated only for the norm-one units $\hat\cO^1$ of an
$R$-order $\cO$, not for a general open subgroup; the same proof extends, since it uses only that $\hat\cO^1$ is open (proof below); this matters because $U^{(q)}$ need not be the unit group of an order (e.g.\ $\Uone{N}$-levels).
In Lean, (SA$_U$) for all admissible $U$ is registered as the hypothesis $\HypSA$, \lean{Langlandscape.Hyp.SA} (shape
\lean{Langlandscape.Tree.SAat}, quantified over \pathtt{Tree.IsQAdmissible}); this is a registered hypothesis, not a proof, and the
theorem stays \stcited. The derivation of (SA$_U$) from density below is not formalised. Per level, $\HypSA$ is replaced by
Prop.~\ref{prop:s02-BG}.

\begin{proof}[Proof of (SA$_U$) from the theorem]
The inclusion ``$\supseteq$'' is clear. $U^{(q)}$ is open in $\Bhat^{\times,(q)}$, so $U^{1,(q)}=U^{(q)}\cap\Bhat^{1,(q)}$ is open in
$\Bhat^{1,(q)}$ (subspace topology; it agrees with the restricted-product topology, as $\cO_\ell^\times\cap B^1_\ell=\cO^1_\ell$ and
$U^1_\ell=\cO^1_\ell$ for almost all $\ell$). Let $y\in\Bhat^1$ and write $y=y_qy^{(q)}$. Then $y^{(q)}\in\Bhat^{1,(q)}$, and
$y^{(q)}U^{1,(q)}$ is an open neighbourhood of $y^{(q)}$. By density there is $\beta\in B^1$ with $\beta^{(q)}\in y^{(q)}U^{1,(q)}$, i.e.\
$u:=(\beta^{(q)})^{-1}y^{(q)}\in U^{1,(q)}$. Put $h:=\beta_q^{-1}y_q$; it lies in $B^1_q$ because $\Nrd\beta=1$ and $\Nrd y_q=1$. Then,
since ideles supported at disjoint sets commute,
\[\beta\cdot h\cdot u=\big(\beta_q\beta_q^{-1}y_q\big)\big(\beta^{(q)}(\beta^{(q)})^{-1}y^{(q)}\big)=y_qy^{(q)}=y.\qedhere\]
\end{proof}

\begin{definition}[bipartite Brandt graph]\label{def:s02-BG}
For $q$-admissible $U$, $\BG(U,q)$ is the finite bipartite multigraph with $V_0=V_1=\XB{U}$ (we write $(x,0)$ and $(x,1)$ for
$x\in\XB U$), edge set $E=\XB{U'}$, $U'=\Uzero{q}U$, and
\[o[y]=(s_0[y],0),\quad t[y]=(s_1[y],1),\qquad s_0[y]=[y],\qquad s_1[y]=[y\alpha]\qquad(\XB{U'}\to\XB{U}).\]
The weights are $w_{(x,j)}=w_x$ (level $U$) and $w_{[y]}=w_y$ (level $U'$), with $w_x=|\Gamma_x/(\Gamma_x\cap\{\pm1\})|$ as in
Def.~\ref{def:s01-classset}. This is the graph that HACS (and the Python track) computes.
\end{definition}
\noindent\stproved{} (well-definedness and finiteness). \emph{Proof.} $s_0$ is well defined since $U'\subseteq U$. For $s_1$: if
$y'=\gamma yu$ with $\gamma\in B^\times$, $u\in U'$, then $y'\alpha=\gamma y\alpha\,(\alpha^{-1}u\alpha)$, and $\alpha^{-1}u\alpha\in U$:
at $q$ because $\alpha^{-1}\Iw_q\alpha\subseteq\GL_2(\Zq)$ (Lemma~\ref{lem:s02-tree}(d)), away from $q$ because $\alpha$ is
supported at $q$. Finiteness: $\XB U$ and $\XB{U'}$ are finite (Prop.~\ref{prop:s01-finite}). Note also that $s_1$ is surjective,
$[x]=s_1[x\alpha^{-1}]$, and so is $s_0$, $[x]=s_0[x]$. \hfill$\square$

\begin{sloppypar}\noindent\emph{Lean status.} \stformal{} (graph, well-definedness, finiteness):
\lean{Langlandscape.Tree.BG}, \lean{Langlandscape.Tree.BGConnected}, \lean{Langlandscape.Tree.s0}, \lean{Langlandscape.Tree.s1},
\lean{Langlandscape.Tree.s0\_mk}, \lean{Langlandscape.Tree.s1\_mk}, \lean{Langlandscape.Tree.alphaElt\_inv\_conj\_mem}
($\alpha^{-1}U'\alpha\subseteq U$), \lean{Langlandscape.Tree.finite\_XB\_of\_isQLevel}, \lean{Langlandscape.Tree.finite\_XB\_U0q},
\lean{Langlandscape.Tree.relIndex\_U0q} ($[U:U']=q+1$). In Lean $s_1$ is the S01 translation map
$\XB{U'}\to\XB U$, $[y]\mapsto[y\alpha]$ (\pathtt{Brandt.translateMap}); the graph is defined for every $q$-level
(\texttt{IsQLevel}, no norm condition), and the weights are not part of the Lean structure.\end{sloppypar}

\begin{proposition}[SA for one level as a finite check]\label{prop:s02-BG}
Let $U$ be $q$-admissible. The maps
\begin{gather*}
\phi_0:\GU\backslash V_0(\cT_q)\to\XB U,\ \GU gv_0\mapsto[g];\\
\phi_1:\GU\backslash V_1(\cT_q)\to\XB U,\ \GU g\alpha v_0\mapsto[g\alpha];\\
\phi_E:\GU\backslash E(\cT_q)\to\XB{U'},\ \GU ge_0\mapsto[g]
\end{gather*}
($g\in B^1_q$) are well defined and \emph{injective}, and $\phi=(\phi_0,\phi_1,\phi_E):\GU\backslash\cT_q\to\BG(U,q)$ is a locally
surjective morphism of graphs. The following are equivalent:
\begin{enumerate}[label=(\roman*)]
\item $\phi_0$ is surjective;
\item (SA$_U$) holds and $\Xone(U)\to\XB{U}$ is surjective;
\item $\BG(U,q)$ is connected.
\end{enumerate}
Moreover:
\begin{enumerate}[label=(\alph*)]
\item (i) implies that $\phi_0,\phi_1,\phi_E$ are all bijective, so $\phi$ is an isomorphism of graphs $\GU\backslash\cT_q\cong\BG(U,q)$.
  This implication does not use the connectedness of $\cT_q$.
\item In all cases, $\phi$ is an isomorphism of $\GU\backslash\cT_q$ onto a union of connected components of $\BG(U,q)$ containing
  $([1],0)$; using that $\cT_q$ is connected, onto the connected component of $([1],0)$.
\end{enumerate}
\end{proposition}
\begin{sloppypar}\noindent\emph{Status.} \stproved{} (new in this plan). No strong approximation is used. The connectedness of $\cT_q$
(Lemma~\ref{lem:s02-tree}(a), \stcited) enters only in (i)$\Rightarrow$(iii) and in the last clause of (b); the implications
(iii)$\Rightarrow$(i)$\Leftrightarrow$(ii), the injectivity, local surjectivity, and (a) are tree-free.
In particular Thm.~\ref{thm:s02-SA} together with Hasse--Schilling (Lemma~\ref{lem:s02-normone}(b)) implies (ii), hence (i) and the
bijections of (a) (tree-free), and hence (iii) (using connectedness of $\cT_q$). As hypotheses, SA alone does not give (ii).
Conversely, for an explicit $(B,U,q)$ the condition (iii) is a finite computation, and it discharges (SA$_U$) and the surjectivity
of $\Xone(U)\to\XB U$ for that level. Classically, SA enters at this point through the definite analogue of Ihara's lemma:
\lit{diamond-taylor\_1994\_non-optimal-levels-of-mod-l-modular-representations}{Lemma~2 and Cor.~1, pp.~445--446} show that $T_\ell-2$, $S_\ell-1$ annihilate $A[l]$, $A$ the kernel of
$(f_1,f_2)\mapsto f_1+\varphi_pf_2$ on $S_k(U;\Q_l/\Z_l)^2$, for almost all $\ell\equiv1\bmod d$ (Lemma~2; hypotheses $p\nmid N_U l$), whence the Eisenstein statement (Cor.~1), and the paper says this lemma ``plays the role of Ihara's result'' (p.~445).
\emph{Lean status.} \stformal{} (all tree-free parts; in the coset model of Lemma~\ref{lem:s02-tree}, with $\phi_0,\phi_1$ combined
into one map \texttt{phiV} on all vertex orbits to $\XB U\sqcup\XB U$):
\lean{Langlandscape.Tree.phiV\_injective}, \lean{Langlandscape.Tree.phiE\_injective},
\lean{Langlandscape.Tree.phiV\_oO}, \lean{Langlandscape.Tree.phiV\_tO} (Step~3),
\lean{Langlandscape.Tree.phi\_locSurj\_type0}, \lean{Langlandscape.Tree.phi\_locSurj\_type1} (Step~4);
(i) is \lean{Langlandscape.Tree.surjType0} (every class of $\XB U$ is $[g]$ with $g\in\SL_2(\Qq)$), and
\lean{Langlandscape.Tree.surjType0\_iff} identifies it with the surjectivity of $\phi_0$;
\lean{Langlandscape.Tree.surjType0\_of\_BGConnected} ((iii)$\Rightarrow$(i)),
\lean{Langlandscape.Tree.bijective\_of\_BGConnected} ((iii)$\Rightarrow$ bijections),
\lean{Langlandscape.Tree.bijective\_of\_surjType0} ((a)), \lean{Langlandscape.Tree.surj\_of\_SA} ((ii)$\Rightarrow$(i)),
\lean{Langlandscape.Tree.SA\_of\_surj} ((i)$\Rightarrow$(ii)). None of these takes a \texttt{Hyp.*} argument.
(i)$\Rightarrow$(iii) is \lean{Langlandscape.Tree.BGConnected\_of\_surj}, formalized modulo \texttt{Hyp.BTtree}; the ``in particular''
is \lean{Langlandscape.Tree.BGConnected\_of\_SA\_HS}, formalized modulo $\HypSA$, $\HypHS$, \texttt{Hyp.BTtree} (and
\lean{Langlandscape.Tree.surjType0\_of\_SA\_HS} for (i), modulo $\HypSA$, $\HypHS$ only). The last clause of (b) (image $=$ the
component of $([1],0)$) also uses connectedness of $\cT_q$; it has no separate Lean declaration (not formalised; it would take
\texttt{Hyp.BTtree}).
\emph{Hypothesis on $U$ in Lean} (audit note): all of these except \pathtt{BGConnected\_of\_SA\_HS} and \pathtt{surjType0\_of\_SA\_HS}
assume only that $U$ is a $q$-level (\texttt{IsQLevel}: $U\le\hat\cO^\times$ of finite index, factoring at $q$), not that $U$ is
$q$-admissible; the condition $\Nrd U=\Zhat^\times$ is not needed for the proposition, in line with the proof above (Step~2 uses only
$\Nrd U\subseteq\Zhat^\times$).
\emph{Deviation from the plan tex (refinement).} The tex's ``Moreover, in all cases the maps are injective and $\phi$ is an isomorphism
onto the component of $(1,0)$'' is split here into the injectivity/local surjectivity in the statement, (a) ((i)$\Rightarrow$ bijections,
tree-free) and (b) (union of components, tree-free; single component, with the tree).\end{sloppypar}

\begin{proof}
We treat the three kinds of cells uniformly. For $c\in\{v_0,v_1,e_0\}$ let
\[a_{v_0}=1,\ a_{v_1}=\alpha,\ a_{e_0}=1;\qquad U_{v_0}=U_{v_1}=U,\ U_{e_0}=U'.\]
Then $S_c=\SL_2(\Qq)\cap a_cU_{c,q}a_c^{-1}$ in all three cases (the list above), every cell of the kind of $c$ is $gc$ with
$g\in B^1_q$, and $\phi_c(\GU gc)=[ga_c]\in\XB{U_c}$. (For $c=v_1$: $g\alpha v_0=gv_1$.)

\emph{Step 0: an identity.} Let $h\in B_q^\times$, $\gamma\in\GU$ and $W\in\{U,U'\}$. Then $[\gamma_qh]=[h]$ in $\XB W$. Indeed
$\gamma_qh=\gamma\cdot h\cdot(\gamma^{(q)})^{-1}$ (the factors $h$ and $\gamma^{(q)}$ commute), and $(\gamma^{(q)})^{-1}\in U^{(q)}=W^{(q)}$
because $\gamma^{(q)}\in U^{(q)}$ by the definition of $\GU$ (no product decomposition of $U^{(q)}$ is used).

\emph{Step 1: $\phi_c$ is well defined.} If $gc=g'c$ with $g,g'\in B^1_q$, then $g^{-1}g'\in S_c\subseteq a_cU_{c,q}a_c^{-1}$, so
$g'a_c=ga_c\,(a_c^{-1}g^{-1}g'a_c)\in ga_cU_{c,q}$ and $[g'a_c]=[ga_c]$. If $g'c=\gamma_qgc$ with $\gamma\in\GU$, the same argument
and Step~0 (with $h=ga_c$) give $[g'a_c]=[\gamma_qga_c]=[ga_c]$.

\emph{Step 2: $\phi_c$ is injective (the norm argument).} Let $g,g'\in B^1_q$ with $[g'a_c]=[ga_c]$ in $\XB{U_c}$: $g'a_c=\beta ga_cu$
with $\beta\in B^\times$, $u\in U_c$. Taking reduced norms (ideles, commutative) and using $\Nrd g=\Nrd g'=1$ gives
$\Nrd a_c=\Nrd\beta\,\Nrd a_c\,\Nrd u$, so $\Nrd\beta=\Nrd u^{-1}$. Here $\Nrd\beta\in\Q_{>0}$ ($B$ is definite) and
$\Nrd u\in\Zhat^\times$ ($U_c\subseteq\hat\cO^\times$). As $\Q_{>0}\cap\Zhat^\times=\{1\}$, we get $\Nrd\beta=1$ and $\Nrd u=1$, so
$\Nrd u_q=1$. Away from $q$ the identity reads $1=\beta^{(q)}u^{(q)}$, so $\beta^{(q)}=(u^{(q)})^{-1}\in U^{(q)}$. Hence
$\beta\in\GU$. At $q$ it reads $g'=\beta_qg\,(a_cu_qa_c^{-1})$ with $a_cu_qa_c^{-1}\in\SL_2(\Qq)\cap a_cU_{c,q}a_c^{-1}=S_c$, so
$g'c=\beta_q\,gc$, i.e.\ $\GU g'c=\GU gc$. (This is Lemma~\ref{lem:s02-normone}(a) in the case at hand; only
$\Nrd(B^\times)\subseteq\Q_{>0}$ and $\Nrd U_c\subseteq\Zhat^\times$ were used.)

\emph{Step 3: $\phi$ is a graph morphism.} Elements of $\SL_2(\Qq)$ preserve types, so for $g\in B^1_q$ the edge $ge_0$ has
$o(ge_0)=gv_0$ and $t(ge_0)=gv_1=g\alpha v_0$. Hence $\phi_0(\bar o(\GU ge_0))=[g]=s_0[g]$ and
$\phi_1(\bar t(\GU ge_0))=[g\alpha]=s_1[g]$, i.e.\ $\phi\circ\bar o=o\circ\phi_E$ and $\phi\circ\bar t=t\circ\phi_E$.

\emph{Step 4: local surjectivity.} \emph{Type $0$.} Let $g\in B^1_q$; the vertex $\GU gv_0$ maps to $([g],0)$. Let $[y]\in\XB{U'}$ be an
edge of $\BG$ at $([g],0)$, i.e.\ $[y]=[g]$ in $\XB U$: $y=\beta gu$ with $\beta\in B^\times$, $u\in U$. In $\XB{U'}$ we have
$[y]=[gu]=[gu_q]$, since $u^{(q)}\in U'$ and $u^{(q)}$ commutes with $gu_q$. Put $d:=\iota_q^{-1}\mathrm{diag}(\det\iota_q(u_q),1)$ and
$k:=u_qd^{-1}$. Then $d\in U'_q$ (lower left entry $0$) and $k\in U^1_q$, so $[y]=[gk]$ with $gk\in B^1_q$. Thus
$[y]=\phi_E(\GU gke_0)$, and $o(gke_0)=gkv_0=gv_0$.
\emph{Type $1$.} Let $g\in B^1_q$; the vertex $\GU g\alpha v_0$ maps to $([g\alpha],1)$. Let $[y]$ be an edge with $s_1[y]=[g\alpha]$:
$y\alpha=\beta g\alpha u$ with $\beta\in B^\times$, $u\in U$. Then $y=\beta g\,(\alpha u_q\alpha^{-1})\,u^{(q)}$, so $[y]=[g\alpha u_q\alpha^{-1}]$ in
$\XB{U'}$. With $d,k$ as before ($u_q=kd$), $\alpha u_q\alpha^{-1}=(\alpha k\alpha^{-1})\,d$, because the diagonal $d$ commutes with $\alpha$, and
$d\in U'_q$. So $[y]=[h]$ with $h:=g\alpha k\alpha^{-1}\in B^1_q$, i.e.\ $[y]=\phi_E(\GU he_0)$, and
$t(he_0)=h\alpha v_0=g\alpha kv_0=g\alpha v_0$.
So for every vertex $w$ of $\GU\backslash\cT_q$, $\phi$ maps the edges at $w$ onto the edges of $\BG(U,q)$ at $\phi(w)$. (An edge of
$\GU\backslash\cT_q$ at $\GU P$ is an orbit $\GU e$ with $\bar o(\GU e)=\GU P$, resp.\ $\bar t$; the edges $gke_0$, $he_0$ above have
an end equal to the given representative.)

\emph{Step 5: a graph lemma.} Let $\psi:G\to H$ be a locally surjective morphism of bipartite multigraphs, and let
$\psi(G)\subseteq H$ be its image (image vertices and image edges). (1) If an edge $\varepsilon$ of $H$ has an end in $\psi(V(G))$, then
$\varepsilon\in\psi(E(G))$, and both ends of $\varepsilon$ lie in $\psi(V(G))$. Indeed, if, say, $o(\varepsilon)=\psi(w)$, local
surjectivity gives $e$ at $w$ with $\psi(e)=\varepsilon$, and $t(\varepsilon)=\psi(t(e))$. So $\psi(G)$ is a union of connected
components of $H$. (2) If $\psi$ is injective on vertices and edges, it is an isomorphism $G\to\psi(G)$: it is bijective onto
$\psi(G)$, and the inverse commutes with $o,t$ because $\psi$ does. (3) If $G$ is connected, so is $\psi(G)$ (a morphism maps the
generating relation $o(e)\sim t(e)$ to $o(\psi e)\sim t(\psi e)$), so $\psi(G)$ is a single connected component.

\emph{Proof of (b).} By Steps 2--5(1),(2), $\phi$ is an isomorphism onto $\phi(\GU\backslash\cT_q)$, a union of components; it
contains $\phi_0(\GU v_0)=([1],0)$. \textbf{[tree]} The tree $\cT_q$ is connected (Lemma~\ref{lem:s02-tree}(a)), so the quotient
$\GU\backslash\cT_q$ is connected: the image of a path in $\cT_q$ is a path in the quotient. By Step~5(3), the image is the component
of $([1],0)$.

\emph{(iii)$\Rightarrow$(i)} (tree-free). By Step~5(1) the image of $\phi$ is a non-empty union of components. If $\BG(U,q)$ is
connected, it is everything; in particular $\phi_0$ is surjective.

\emph{(i)$\Rightarrow$ the bijections, i.e.\ (a)} (tree-free). Assume (i). Every vertex $(x,1)$ of $\BG$ is the end $t[x\alpha^{-1}]$ of
an edge, whose other end $(s_0[x\alpha^{-1}],0)$ lies in the image by (i). By Step~5(1) this edge and $(x,1)$ lie in the image, so
$\phi_1$ is surjective. Every edge $[y]$ has the end $(s_0[y],0)$ in the image, so it lies in the image: $\phi_E$ is surjective.
With Step~2 all three maps are bijective, and by Step~5(2) $\phi$ is an isomorphism of graphs. No connectedness was used.

\emph{(i)$\Rightarrow$(iii).} By (a), $\BG(U,q)\cong\GU\backslash\cT_q$. \textbf{[tree]} This is the only place where connectedness of
$\cT_q$ enters: the quotient of the connected graph $\cT_q$ is connected (as in the proof of (b)). Hence $\BG(U,q)$ is connected.
Without connectedness of $\cT_q$ one only gets that $\BG(U,q)$ is isomorphic to $\GU\backslash\cT_q$.

\emph{(i)$\Rightarrow$(ii)} (tree-free). Let $y\in\Bhat^1$. By (i) there is $g\in B^1_q$ with $[y]=[g]$ in $\XB U$: $y=\beta gu$,
$\beta\in B^\times$, $u\in U$. As in Step~2, $\Nrd\beta=\Nrd u^{-1}\in\Q_{>0}\cap\Zhat^\times=\{1\}$, so $\beta\in B^1$ and $u\in U^1=U^1_q\times
U^{1,(q)}$. Then $y=\beta\cdot(gu_q)\cdot u^{(q)}\in B^1(B^1_q\times U^{1,(q)})$, which is (SA$_U$). Surjectivity of $\Xone(U)\to\XB U$:
every class is $[g]$ with $g\in B^1_q\subseteq\Bhat^1$ by (i), and $[g]$ is the image of the class of $g$ in $\Xone(U)$.

\emph{(ii)$\Rightarrow$(i)} (tree-free; the seed's argument). Let $[x]\in\XB U$. By surjectivity of $\Xone(U)\to\XB U$ there is
$y\in\Bhat^1$ with $[x]=[y]$. By (SA$_U$), $y=\beta hu$ with $\beta\in B^1$, $h\in B^1_q$, $u\in U^{1,(q)}\subseteq U$. Hence
$[x]=[y]=[h]=\phi_0(\GU hv_0)$.

\emph{The ``in particular''.} Thm.~\ref{thm:s02-SA} gives (SA$_U$). Since $\Nrd U=\Zhat^\times$ ($U$ is $q$-admissible), Hasse--Schilling
gives the surjectivity of $\Xone(U)\to\XB U$ (Lemma~\ref{lem:s02-normone}(b)). So (ii) holds, and the rest follows from the
implications just proved.
\end{proof}

\begin{proposition}[quotient graph; seed Prop.~quot]\label{prop:s02-quot}
Let $U$ be $q$-admissible and assume (iii) of Prop.~\ref{prop:s02-BG}. By Prop.~\ref{prop:s02-BG} this holds given
Thm.~\ref{thm:s02-SA}, Hasse--Schilling \emph{and} the connectedness of $\cT_q$ (Lemma~\ref{lem:s02-tree}(a), \stcited; \textbf{[tree]}):
$\HypSA\wedge\HypHS\wedge\mathrm{Hyp.BTtree}\Rightarrow$ (iii) (Lean \pathtt{Tree.BGConnected\_of\_SA\_HS}, which takes \texttt{Hyp.BTtree});
the bijections of (a) need only $\HypSA\wedge\HypHS$. As hypotheses, SA alone does not suffice.
(iii) itself is the working hypothesis here.
\begin{enumerate}[label=(\alph*)]
\item There are bijections
  \begin{gather*}
    \GU\backslash\{\text{type-}0\text{ vertices}\}\ \xrightarrow{\ gv_0\mapsto[g]\ }\ \XB{U},\\
    \GU\backslash\{\text{type-}1\text{ vertices}\}\ \xrightarrow{\ g\alpha v_0\mapsto[g\alpha]\ }\ \XB{U},\\
    \GU\backslash\{\text{edges}\}\ \xrightarrow{\ ge_0\mapsto[g]\ }\ \XB{\Uzero{q}U}
  \end{gather*}
  ($g\in B^1_q$). Under these bijections the endpoint maps $o,t$ become $s_0,s_1$. In other words,
  $\GU\backslash\cT_q\cong\BG(U,q)$ as graphs.
\item For $g\in B^1_q$, $\Stab_{\GU}(gv_0)=\Gamma_g$ (level $U$), $\Stab_{\GU}(g\alpha v_0)=\Gamma_{g\alpha}$
  (level $U$) and $\Stab_{\GU}(ge_0)=\Gamma_g$ (level $\Uzero{q}U$). For another representative
  $x'=\gamma xu$ one has $\Gamma_{x'}=\gamma\Gamma_x\gamma^{-1}$. Each $\Gamma_x$ has norm $1$
  (Lemma~\ref{lem:s01-stab}(a)), so it lies in $\GU$ after conjugation.
\item Let $T$ be any Hecke operator supported away from $q$ ($T_\ell,S_\ell,U_\ell$ for $\ell\mid M$,
  $\langle d\rangle$, $U_p$ with the normalisation $\NUp$ of stage 01, and $\Pi_\ell$ for $\ell\mid D$),
  acting compatibly on levels $U$ and $\Uzero{q}U$. Then $s_0^*,s_1^*,s_{0*},s_{1*}$ commute with $T$ on
  Brandt modules $\Brandt{\cdot}{V}$ and on their homological versions (on $\Brandt{\cdot}{V}$, $s_{j*}$ means the trace
  $\mathrm{tr}_j$ of Def.~\ref{def:s02-Mhom}).
\end{enumerate}
\end{proposition}
\begin{sloppypar}\noindent\emph{Status.} (a): \stproved{} given (iii); \stpmc{} via Thm.~\ref{thm:s02-SA} and Hasse--Schilling (through
(ii)$\Rightarrow$(i)$\Rightarrow$(a), which does not use connectedness of the tree). (b): \stproved; \emph{refinement of the plan tex} (which has ``given (iii)''): the three stabiliser identities
hold for every $g\in B^1_q$ without (iii); (iii) is used only to know that every class of $\XB U$, $\XB{U'}$ has such a tree
representative. (c): \stproved{} (unconditional).

\emph{Lean status.} (a) \stformal{} given (iii): \lean{Langlandscape.Tree.quotVEquiv}, \lean{Langlandscape.Tree.quotEEquiv},
\lean{Langlandscape.Tree.quotEquiv\_oO}, \lean{Langlandscape.Tree.quotEquiv\_tO}; formalized modulo $\HypSA$, $\HypHS$ (no
\texttt{Hyp.BTtree}): \lean{Langlandscape.Tree.quotVEquiv\_of\_SA\_HS}, \lean{Langlandscape.Tree.quotEEquiv\_of\_SA\_HS}.
(b) \stformal{} for the three stabiliser identities, for every $g\in\SL_2(\Qq)$ and without (iii) (as in the refinement above):
\lean{Langlandscape.Tree.stab\_gv0\_eq\_Gam}, \lean{Langlandscape.Tree.stab\_gav0\_eq\_Gam}, \lean{Langlandscape.Tree.stab\_ge0\_eq\_Gam};
the conjugation of $\Gamma_x$ under a change of representative has no separate declaration here.
(c) \stformal{} for \emph{trivial coefficients} $V=R$ and the maps $s_0^*,s_1^*$ (\texttt{s0Pull}, \texttt{s1Pull}) and $s_{0*},s_{1*}$
(\texttt{s0Push}, \texttt{s1Push}): \lean{Langlandscape.Tree.heckeT\_comm\_s0Pull}, \lean{Langlandscape.Tree.heckeT\_comm\_s1Pull},
\lean{Langlandscape.Tree.heckeTHom\_comm\_s0Push}, \lean{Langlandscape.Tree.heckeTHom\_comm\_s1Push}, with the named corollaries
\lean{Langlandscape.Tree.Tl\_comm\_s1Pull}, \lean{Langlandscape.Tree.Sl\_comm\_s1Pull}, \lean{Langlandscape.Tree.diamond\_comm\_s1Pull},
\lean{Langlandscape.Tree.varpi\_comm\_s1Pull} ($U_p$ through \texttt{varpiDiag}), and the level-change half
\lean{Langlandscape.Tree.heckeT\_comm\_levelPull}, \lean{Langlandscape.Tree.heckeTHom\_comm\_levelPush}. \emph{Shape in Lean:} the
operator is \texttt{heckeT U t} for any $t$ with $q$-component $1$ and finite double cosets (\texttt{hfin}); Lean needs no factorisation
of $U$ at $\Sigma$, only the factorisation at $q$ contained in \texttt{IsQLevel}. Not formalised: the traces $\mathrm{tr}_j$ (for
those, Prop.~\ref{prop:s02-Mhomprops}(d) and the adjointness below), general $V$, and $\Pi_\ell$.
\emph{Level change.} $s_1^*$ is factored in Lean as \lean{Langlandscape.Tree.s1Pull\_eq}: $s_1^*=R_\alpha\circ\pi^*$, with $\pi^*$ the
pullback from $U$ to $U'_\alpha=\alpha^{-1}U'\alpha$ (\lean{Langlandscape.Tree.s1Level}, \texttt{levelPull}) and
$R_\alpha$ $=$ \lean{Langlandscape.Brandt.rightTranslateLin}, $(R_\alpha f)[x]=f[x\alpha]$, the right translation from level
$g^{-1}Wg$ to $W$ (here $W=U'$, $g=\alpha$). The commutation of $T$ with $R_\alpha$ is S01's
\pathtt{Brandt.heckeT\_comm\_rightTranslate}.\end{sloppypar}

\begin{proof}
\emph{(a).} (iii)$\Rightarrow$(i) and (i)$\Rightarrow$ bijections, Prop.~\ref{prop:s02-BG}; the identification of $o,t$ with $s_0,s_1$ is
Step~3 there.

\emph{(b).} Let $c\in\{v_0,v_1,e_0\}$ with $a_c,U_c$ as in the proof of Prop.~\ref{prop:s02-BG}, and $g\in B^1_q$, $x:=ga_c$. By
Lemma~\ref{lem:s02-tree}(c), $\Stab_{\SL_2(\Qq)}(gc)=gS_cg^{-1}=\SL_2(\Qq)\cap x\,U_{c,q}\,x^{-1}$. Let $\gamma\in\GU$ fix $gc$. Then
$\gamma\in B^\times$, $\gamma^{(q)}\in U^{(q)}=x^{(q)}U_c^{(q)}(x^{(q)})^{-1}$ (as $x^{(q)}=1$ and $U_c^{(q)}=U^{(q)}$) and
$\gamma_q\in x_qU_{c,q}x_q^{-1}$; so $\gamma\in B^\times\cap xU_cx^{-1}=\Gamma_x$ (level $U_c$). Conversely let $\gamma\in\Gamma_x$. By
Lemma~\ref{lem:s01-stab}(a) ($\Nrd U_c=\Zhat^\times$ by Def.~\ref{def:s02-GU}; only $\subseteq$ is needed) $\Nrd\gamma=1$; $\gamma^{(q)}\in U_c^{(q)}=U^{(q)}$, so $\gamma\in\GU$; and $\gamma_q\in\SL_2(\Qq)\cap x_qU_{c,q}x_q^{-1}$ fixes $gc$. For $c=v_0,v_1,e_0$ this is the
three claimed identities ($x=g$, $g\alpha$, $g$).
For $x'=\gamma xu$ ($\gamma\in B^\times$, $u\in U_c$): $x'U_cx'^{-1}=\gamma xU_cx^{-1}\gamma^{-1}$, so
$\Gamma_{x'}=B^\times\cap\gamma xU_cx^{-1}\gamma^{-1}=\gamma\Gamma_x\gamma^{-1}$. Finally, given (iii), by (a) every class of $\XB{U_c}$ has a
representative $x=ga_c$ with $g\in B^1_q$; if $x'$ is any representative, $x'=\gamma xu$ and $\gamma^{-1}\Gamma_{x'}\gamma=\Gamma_x
=\Stab_{\GU}(gc)\subseteq\GU$.

\emph{(c).} \emph{The maps.} For compact open $W_1,W_2\subseteq\hat\cO^\times$ and $h\in\Bhat^\times$ with $h_p\in U_p$-monoid as in
Def.~\ref{def:s01-hecke} (here always $h_p\in\Sigma_p$), write $W_1hW_2=\coprod_ih_iW_2$ (finite) and
\[[W_1hW_2]:\Brandt{W_2}{V}\to\Brandt{W_1}{V},\qquad([W_1hW_2]f)(x)=\textstyle\sum_ih_{i,p}\cdot f(xh_i).\]
It is well defined, by the computation in the proof of Def.~\ref{def:s01-hecke} with $(U,U)$ replaced by $(W_1,W_2)$: for $w\in W_1$,
$wh_i=h_{\sigma(i)}w_i$ with $w_i\in W_2$, and the representatives may be changed by $W_2$ on the right. In cohomological form
\begin{gather*}s_0^*=[U'\,1\,U],\quad s_1^*=[U'\,\alpha\,U]:\Brandt{U}{V}\to\Brandt{U'}{V},\\
\mathrm{tr}_0=[U\,1\,U'],\quad \mathrm{tr}_1=[U\,\alpha^{-1}\,U']:\Brandt{U'}{V}\to\Brandt{U}{V}.\end{gather*}
Since $U'\subseteq U$ and $\alpha^{-1}U'\alpha\subseteq U$, $U'1U=U$ and $U'\alpha U=\alpha U$ are single cosets, so $s_0^*f=f$ and
$s_1^*f=f(\cdot\,\alpha)$ ($\alpha_p=1$): pullback along $s_0,s_1$. (For $s_1^*$: first restrict $f$ from level $U$ to
$U'_\alpha=\alpha^{-1}U'\alpha\subseteq U$, then right-translate by $\alpha$ from level $\alpha^{-1}U'\alpha$ to $U'$; this is the
factorisation \lean{Langlandscape.Tree.s1Pull\_eq} through \lean{Langlandscape.Brandt.rightTranslateLin}.) The other two are the traces (the plan's ``trace $s_{j*}$'' on $\Brandt{\cdot}{V}$; we follow \S3 and write $\mathrm{tr}_j$, keeping
$s_{j*}$ for the homological pushforward),
$\mathrm{tr}_0f(x)=\sum_{u\in U/U'}f(xu)$ and $\mathrm{tr}_1f(x)=\sum_if(xu_i\alpha^{-1})$ for $U\alpha^{-1}U'=\coprod_iu_i\alpha^{-1}U'$,
i.e.\ $u_i$ running over $U/U'_\alpha$, $U'_\alpha=\alpha^{-1}U'\alpha$; here all coset representatives can be taken in $B_q^\times$,
because $U$ and $U'$ agree away from $q$, so $u_p=u_{i,p}=1$. These are exactly the maps $\mathrm{tr}_0,\mathrm{tr}_1$ of
Def.~\ref{def:s02-Mhom} (whose $u_p$-twist is trivial here). In homological form (the coinvariant modules
$\Brandt{W}{V}_*$, generated by symbols $[x,v]$ with $[\gamma xw,v]=[x,w_pv]$ for $\gamma\in B^\times$, $w\in W$; this is
$\bigoplus_{[x]}V_{x_p^{-1}\Gamma_xx_p}$) the pushforwards are $s_{0*}[y,v]=[y,v]$ and $s_{1*}[y,v]=[y\alpha,v]$; the second is well
defined since $[yw\alpha,v]=[y\alpha\cdot\alpha^{-1}w\alpha,v]=[y\alpha,w_pv]$ for $w\in U'$.

\emph{The Hecke operators.} Let $h$ be supported at a finite set $\Sigma$ of primes with $q\notin\Sigma$. \emph{Hypothesis} (part of
``acting compatibly''; it is not part of $q$-admissibility, which only gives $U=U_q\times U^{(q)}$): $U=U_\Sigma\times U^\Sigma$ with
$U_\Sigma=\prod_{v\in\Sigma}U_v$. It holds for all standard levels, which are factorisable, $U=\prod_vU_v$ (Def.~\ref{def:s01-levels}).
Then $UhU=U_\Sigma h_\Sigma U_\Sigma\times U^\Sigma$, written below for factorisable $U$ as
$UhU=\prod_{v\in\Sigma}U_vh_vU_v\times\prod_{v\notin\Sigma}U_v$, so $UhU=\coprod_ih_iU$ with all $h_i$ supported at $\Sigma$. As
$U'_v=U_v$ for $v\in\Sigma$ and $U'_q=\Iw_q$, also $U'hU'=\coprod_ih_iU'$ with the \emph{same} $h_i$. This is what ``acting compatibly''
means: the operator is given at both levels by the same data $(h_i,A_i)$, namely $(Tf)(x)=\sum_iA_if(xh_i)$ with $A_i=h_{i,p}\cdot$
(cohomological), resp.\ $T[x,v]=\sum_i[xh_i,A_iv]$ for the homological normalisation, for endomorphisms $A_i$ of $V$ satisfying the compatibility $A_{\sigma(i)}=u_{i,p}A_iu_p^{-1}$ (where $uh_i=h_{\sigma(i)}u_i$; e.g.\ $A_i=1$ when $h_p=1$), which is needed for well-definedness; the homological $U_p$ is not spelled out in this chapter
(\pathtt{if:s02-indef}: transposed right-translation operators). All operators listed in (c) are of this form: $T_\ell=T_{\varpi_\ell}$,
$S_\ell=T_{z_\ell}$, $U_\ell=T_{\varpi_\ell}$ ($\ell\mid M$), $\langle d\rangle=T_{\delta_d}$ with $\delta_d$ supported at the primes dividing
the level away from $q$, $U_p=T_{\varpi_p}$ with representatives $\varpi_j$ supported at $p$ (Def.~\ref{def:s01-NUp}), and $\Pi_\ell$,
$\ell\mid D$; none is supported at $q$, because $q\nmid DMp$ and $q$ does not divide the level of $U$ (as $U_q=\GL_2(\Zq)$).

\emph{Commutation.} Let $g_j\in B_q^\times$ be the representatives of one of the four maps $[W_1gW_2]$ above (with $W_1,W_2\in\{U,U'\}$,
$g\in\{1,\alpha,\alpha^{-1}\}$); $g_{j,p}=1$. Since $g_j$ and $h_i$ are supported at the disjoint sets $\{q\}$ and $\Sigma$, $g_jh_i=h_ig_j$.
Cohomologically, for $f\in\Brandt{W_2}{V}$:
\[(T[W_1gW_2]f)(x)=\sum_iA_i\sum_jf(xh_ig_j)=\sum_j\sum_iA_if(xg_jh_i)=([W_1gW_2]Tf)(x).\]
Homologically, $Ts_{1*}[y,v]=\sum_i[y\alpha h_i,A_iv]=\sum_i[yh_i\alpha,A_iv]=s_{1*}T[y,v]$, and $s_{0*}$ is the identity on symbols.
This is Prop.~\ref{prop:s01-relations}(c) for $R_\alpha$ and $\rho_\alpha$, extended to the traces.
\end{proof}

\begin{remark}[stabilisers modulo $\pm1$; index bookkeeping]\label{rem:s02-quot-index}
\emph{Addition to the plan tex} (new label; it does not restate the index factors of the norm maps, which are
Def.~\ref{def:s02-Mhom}). Assume (iii). (1) $\GU\cap\{\pm1\}=\{\pm1\}\cap U=\Gamma_x\cap\{\pm1\}$ for every cell representative $x$
(the constant $\varepsilon$ of Def.~\ref{def:s02-Mhom}). It acts trivially on $\cT_q$, and it is exactly the kernel of the action: if $g\in\SL_2(\Qq)$ fixes every
vertex, then $g\in S_{v_0}=\SL_2(\Zq)$; for every primitive $w\in\Zq^2$ and $n\ge1$ the lattice $L=\Zq w+q^n\Zq^2$ satisfies
$gL=\lambda L$ with $[\Zq^2:gL]=[\Zq^2:L]$, so $\lambda\in\Zq^\times$ and $gL=L$, i.e.\ $gw\in\Zq w+q^n\Zq^2$; applied to
$e_1,e_2,e_1+e_2$ this gives $g\equiv\lambda_nI\bmod q^n$ for all $n$, so $g=\lambda I$ is scalar, and $\det g=1$ gives $\lambda=\pm1$. Hence the effective stabilisers are
$\Gamma_x/(\Gamma_x\cap\{\pm1\})$, of order $w_x$, and $|\Gamma_x|/|\Gamma_y|=w_x/w_y$ for cells at the two levels.
(2) Inclusions: for $g\in B^1_q$, $\Gamma_g(U')\subseteq\Gamma_g(U)$ and $\Gamma_g(U')\subseteq\Gamma_{g\alpha}(U)$, the latter because
$\alpha^{-1}\Iw_q\alpha\subseteq\GL_2(\Zq)$; these are $\Gamma_e\hookrightarrow\Gamma_{o(e)}$, $\Gamma_e\hookrightarrow\Gamma_{t(e)}$ for $e=ge_0$.
(3) Fibre counts: for every $[x]\in\XB U$ and $j\in\{0,1\}$,
\[\sum_{[y]\in s_j^{-1}[x]}\frac{|\Gamma_x|}{|\Gamma_y|}=q+1,\qquad\text{hence}\qquad\sum_{[y]\in\XB{U'}}\frac1{|\Gamma_y|}=(q+1)\sum_{[x]\in\XB U}\frac1{|\Gamma_x|}.\]
Proof: take the tree representative $P$ of the vertex $(x,j)$. The edges at $P$ are $gke_0$ with $k$ running over
$S_P/(S_P\cap\Stab e)$, which for $P=v_0$ is $\SL_2(\Zq)/(\SL_2(\Zq)\cap\Iw_q)\cong\SL_2(\F_q)/B(\F_q)$, of order $q+1$ (the same for
$v_1$ after conjugating by $\alpha$). $\Gamma_P$-orbits on the edges at $P$ correspond bijectively to $\GU$-orbits of edges at $\GU P$,
i.e.\ by (a) to $s_j^{-1}[x]$, and the orbit of $e$ has $|\Gamma_P|/|\Gamma_e|$ elements. Sizes of $\Gamma$ do not depend on representatives by
(b). (The mass identity is also the level-change formula \texttt{if:s01-mass} with $[U:U']=q+1$, unconditionally;
$[U:U']=q+1$ is \lean{Langlandscape.Tree.relIndex\_U0q}.) \stproved{}
given (iii); these are the checks H02.2(b).
\end{remark}

\begin{remark}[why norm one; seed Remark~normone]\label{rem:s02-normone}
Let $\Gamma^{\rm full}:=\{\gamma\in B^\times:\gamma^{(q)}\in U^{(q)}\}$, acting on $\cT_q$ through $\iota_q(\gamma_q)\in\GL_2(\Qq)$.
\begin{enumerate}[label=(\arabic*)]
\item $\Nrd\gamma\in\Q_{>0}\cap\Z_\ell^\times$ for all $\ell\ne q$, so $\Nrd(\Gamma^{\rm full})\subseteq q^{\Z}$. Since
  $\tau(gP)=\tau(P)+\ord_q\det g\bmod2$ (Def.~\ref{def:s02-tree}), the elements of odd norm valuation exchange the types; the centre
  $\Qq^\times$ of $\GL_2(\Qq)$ acts trivially.
\item The map $\Gamma^{\rm full}\backslash\{\text{all vertices}\}\to B^\times\backslash\Bhat^\times/\Qq^\times U$, $hv_0\mapsto[h]$
  ($h\in B_q^\times$), is well defined and injective, and it is bijective if (i) of Prop.~\ref{prop:s02-BG} holds.
\item $B^\times\backslash\Bhat^\times/\Qq^\times U=\XB U/\langle z_q\rangle$, the orbits of right multiplication by the central idele $z_q$
  ($q$ at $q$, $1$ elsewhere). Modulo $\Q^\times$, $z_q$ equals the scalar idele $y:=(q^{-1})^{(q)}$ ($1$ at $q$, $q^{-1}$ at $\ell\ne q$).
  If $y\in U$ (e.g.\ $U=\Uzero{N}$, or any $U\supseteq\Zhat^{\times,(q)}$) it acts trivially. At $U=\Uone{N}$, $q\nmid N$, it acts as
  $[x]\mapsto[x\delta_{q^{-1}}]$, the diamond $\langle q\rangle^{-1}$ at the primes dividing the level (Def.~\ref{def:s01-hecke}).
\end{enumerate}
So, under (i), the full group gives \emph{one} copy of $\XB{U}/\langle q\rangle$, mixing the two types. At maximal level this is the classical picture
with $h$ vertices \lit{vigneras\_1980\_arithmetique-des-algebres-de-quaternions}{Ch.~V, \S3.B, Prop.~3.3, p.~147: the number of vertices
of the quotient graph is the class number; her group is $\cO[1/p]^\times/\Z[1/p]^\times$ and her vertex orders are stabilisers modulo $\pm1$}; see also \lit{voight\_2021\_quaternion-algebras}{41.2.13, p.~769 and Example 28.5.19,
p.~487} (both page-checked). The norm-one group gives two copies of $\XB{U}$ (Prop.~\ref{prop:s02-quot}(a)), which is what the Ribet
sequence needs.
\stproved{} (1), (3) and the injectivity in (2); the surjectivity in (2) \stproved{} given (i), \stpmc{} via
Thm.~\ref{thm:s02-SA} and Hasse--Schilling.
\emph{Deviation from the plan tex.} The tex states the bijection in (2) unconditionally and marks the remark \stsketch; here
injectivity is proved unconditionally, bijectivity under (i), and the status is upgraded accordingly.
The seed wording ``the central element $q$ acts trivially'' is imprecise: $q\in\Gamma^{\rm full}$ only if $q\in U_\ell$ for all $\ell\ne q$,
which fails at $\Uone{N}$-levels ($N\ge3$ with $q\not\equiv1\bmod N$). What is true is that the centre of $\GL_2(\Qq)$ acts trivially
on $\cT_q$.
\end{remark}
\begin{proof}
(1) $\Nrd\gamma>0$ as $B$ is definite, and $\Nrd\gamma=\Nrd\gamma_\ell\in\Nrd U_\ell\subseteq\Z_\ell^\times$ for $\ell\ne q$; a positive
rational that is an $\ell$-unit for all $\ell\ne q$ is in $q^\Z$. If $P=[a\Zq^2]$, then $gP=[ga\Zq^2]$ and $\tau(gP)\equiv\ord_q\det g+
\ord_q\det a$; for $\gamma\in\Gamma^{\rm full}$, $\ord_q\det\iota_q(\gamma_q)=\ord_q\Nrd\gamma$. Scalars $\lambda$ send $[L]$ to $[\lambda L]=[L]$.

(2) $\GL_2(\Qq)$ acts transitively on all vertices ($[L]=[a\Zq^2]$ for any basis matrix $a$ of $L$), and the stabiliser of $v_0$ is
$\Qq^\times\GL_2(\Zq)$: $a\Zq^2=\lambda\Zq^2$ iff $\lambda^{-1}a\in\GL_2(\Zq)$. As $\Qq^\times U_q\subseteq\Qq^\times U$, the map
$hv_0\mapsto[h]$ is well defined on vertices. It is $\Gamma^{\rm full}$-invariant by Step~0 of the proof of Prop.~\ref{prop:s02-BG}
(which only used $\gamma^{(q)}\in U^{(q)}$). Injectivity: if $h'=\beta hzu$ with $\beta\in B^\times$, $z\in\Qq^\times$, $u\in U$,
then away from $q$, $1=\beta^{(q)}u^{(q)}$, so $\beta^{(q)}\in U^{(q)}$ and $\beta\in\Gamma^{\rm full}$; at $q$, $h'=\beta_qh(zu_q)$ with $zu_q\in\Stab(v_0)$, so
$h'v_0=\beta_qhv_0$. No norm argument is needed here. Surjectivity under (i): every class of $\XB U$, and hence of the coarser
quotient, is $[g]$ with $g\in B^1_q\subseteq B_q^\times$.

(3) $\Qq^\times=q^\Z\Zq^\times$ and $\Zq^\times\subseteq U_q$, so the double coset space is $\XB U$ modulo $[x]\mapsto[xz_q]$, which is
well defined because $z_q$ is central. With $q\in\Q^\times\subseteq B^\times$ central, $z_q=q\,y$, so $[xz_q]=[qxy]=[xy]$. If $y\in U$ this
is $[x]$. At $U=\Uone N$: at $\ell\mid N$, $y_\ell=q^{-1}I=\mathrm{diag}(1,q^{-1})\cdot\mathrm{diag}(q^{-1},1)$ with
$\mathrm{diag}(q^{-1},1)\in U_{1,\ell}$; at $\ell\nmid Nq$, $y_\ell\in\Z_\ell^\times\subseteq U_\ell$. So $y\in\delta_{q^{-1}}U$ with the diagonal
$\delta_d$ of Def.~\ref{def:s01-hecke}, and $[xy]=[x\delta_{q^{-1}}]$.
\end{proof}
